\subsection{Implementation on CPUs}
\label{sec:hws_impl_cpus}

On the \gls{cpu}, however, this is different since there is no simple library that can be called to gather all the necessary information. 
Therefore, three different command-line utilities are used if and only if they can be found during CMake configuration; otherwise, all related code is ignored, and the respective samples are not gathered. 
To retrieve general \gls{cpu} information, like the core count, socket count, or cache size, \texttt{lscpu}~\cite{qian_lscpu_nodate} is used, for memory-related information, i.e., available, used, and total main memory, the \texttt{free}~\cite{small_free_nodate} utility is used, and for everything else, like frequency or power related information, \texttt{turbostat}~\cite{brown_turbostat_nodate} is used. 
In all three cases, the \textit{subprocess.h} library~\cite{henning_subprocessh_2024} is used to create a new subprocess that, in turn, executes the respective utility, returning the standard output as a string, which is then post-processed, extracting the necessary information inside the \gls{cpu} hardware sampler.  

% lscpu
In case of the \texttt{lscpu} utility, the subprocess simply calls \texttt{lscpu}. 
Note that the cache sizes are only reported as strings instead of integers containing the actual values since the \texttt{lscpu} output is not standardized and varies greatly between \glspl{cpu}. 
%
% free
For the \texttt{free} utility, the command-line flag \texttt{-b} is added to ensure that the memory sizes are output in bytes. 

%turbostat
For \texttt{turbostat}, it is a bit more difficult. 
First, it should be noted that to get the full \texttt{turbostat} information, it has to be executed with root privileges. 
However, in newer \texttt{turbostat} versions, it is possible to get a subset of all possible information even if \texttt{turbostat} is not called with root privileges. 
The hws library handles this distinction during CMake configuration. 
If \texttt{turbostat} can be used with root privileges, the command \code{sudo turbostat -n 1 -i 0.001 -S -q} is executed; if root privileges are not granted, \code{sudo} is omitted. 
This command samples all available information immediately (\code{-i 0.001}) exactly once (\code{-n 1}) and outputs the information in a short summary format (\code{-S}), skipping potential header information (\code{-q}). 
An example \texttt{turbostat} output with the mentioned command is displayed in \autoref{code:turbostat_output}. 
\begin{listing}
    \begin{moderncode}{text}
Avg_MHz  Busy%  Bzy_MHz  TSC_MHz  IPC   IRQ  POLL  C1   C2   POLL%  C1%   C2%     CorWatt  PkgWatt
39       1.77   2212     4686     0.92  818  1     156  819  0.01   3.03  110.73  4.58     264.06
    \end{moderncode}
    \caption{%
    An example output of the turbostat command \code{sudo turbostat -n 1 -i 0.001 -S -q} with root privileges enabled on an \gls{amd} EPYC 9274F \gls{cpu}.
    }
    \label{code:turbostat_output}
\end{listing}
Note that the output of \texttt{turbostat} also depends on the used \gls{cpu}. 
For example, the output in \autoref{code:turbostat_output} was gathered on an \gls{amd} EPYC 9274F \gls{cpu} and consists of $14$ table entries. 
The same command executed on an Intel i5-10210U \gls{cpu}, the output consists of $52$ columns, mostly adding more idle state information but also temperature-related information missing for the \gls{amd} \gls{cpu}.
